\documentclass[a4paper,USenglish,cleveref, autoref, thm-restate]{lipics-v2021}
%This is a template for producing LIPIcs articles.
%See lipics-v2021-authors-guidelines.pdf for further information.
%for A4 paper format use option "a4paper", for US-letter use option "letterpaper"
%for british hyphenation rules use option "UKenglish", for american hyphenation rules use option "USenglish"
%for section-numbered lemmas etc., use "numberwithinsect"
%for enabling cleveref support, use "cleveref"
%for enabling autoref support, use "autoref"
%for anonymousing the authors (e.g. for double-blind review), add "anonymous"
%for enabling thm-restate support, use "thm-restate"
%for enabling a two-column layout for the author/affilation part (only applicable for > 6 authors), use "authorcolumns"
%for producing a PDF according the PDF/A standard, add "pdfa"

%\pdfoutput=1 %uncomment to ensure pdflatex processing (mandatatory e.g. to submit to arXiv)
%\hideLIPIcs  %uncomment to remove references to LIPIcs series (logo, DOI, ...), e.g. when preparing a pre-final version to be uploaded to arXiv or another public repository

%\graphicspath{{./graphics/}}%helpful if your graphic files are in another directory

\bibliographystyle{plainurl}% the mandatory bibstyle

\title{Formalizing the Exponential Blowup in the Transformations between CNF and DNF}

%\titlerunning{} %TODO optional, please use if title is longer than one line

\def\lmu{Ludwig-Maximilians-Universität München}

\author{Leon Raffael Schulz}{\lmu{}, Germany}{Le.Schulz@campus.lmu.de}{https://orcid.org/0009-0007-7666-2034}{}
\author{Martin Desharnais-Schäfer}{\lmu{}, Germany}{desharnais-schaefer@ifi.lmu.de}{https://orcid.org/0000-0002-1830-7532}{}
\author{Jan Johannsen}{\lmu{}, Germany}{jan.johannsen@ifi.lmu.de}{https://orcid.org/0009-0004-2072-4679}{}

\authorrunning{Leon Raffael Schulz and Martin Desharnais-Schäfer and Jan Johannsen}

\Copyright{Leon Raffael Schulz and Martin Desharnais-Schäfer and Jan Johannsen}

\ccsdesc[100]{\textcolor{red}{Replace ccsdesc macro with valid one}} %TODO mandatory: Please choose ACM 2012 classifications from https://dl.acm.org/ccs/ccs_flat.cfm

\keywords{Dummy keyword} %TODO mandatory; please add comma-separated list of keywords

\category{} %optional, e.g. invited paper

\relatedversion{} %optional, e.g. full version hosted on arXiv, HAL, or other respository/website
%\relatedversiondetails[linktext={opt. text shown instead of the URL}, cite=DBLP:books/mk/GrayR93]{Classification (e.g. Full Version, Extended Version, Previous Version}{URL to related version} %linktext and cite are optional

%\supplement{}%optional, e.g. related research data, source code, ... hosted on a repository like zenodo, figshare, GitHub, ...
%\supplementdetails[linktext={opt. text shown instead of the URL}, cite=DBLP:books/mk/GrayR93, subcategory={Description, Subcategory}, swhid={Software Heritage Identifier}]{General Classification (e.g. Software, Dataset, Model, ...)}{URL to related version} %linktext, cite, and subcategory are optional

%\funding{(Optional) general funding statement \dots}%optional, to capture a funding statement, which applies to all authors. Please enter author specific funding statements as fifth argument of the \author macro.

\acknowledgements{}%optional

%\nolinenumbers %uncomment to disable line numbering



%Editor-only macros:: begin (do not touch as author)%%%%%%%%%%%%%%%%%%%%%%%%%%%%%%%%%%
\EventEditors{John Q. Open and Joan R. Access}
\EventNoEds{2}
\EventLongTitle{42nd Conference on Very Important Topics (CVIT 2016)}
\EventShortTitle{CVIT 2016}
\EventAcronym{CVIT}
\EventYear{2016}
\EventDate{December 24--27, 2016}
\EventLocation{Little Whinging, United Kingdom}
\EventLogo{}
\SeriesVolume{42}
\ArticleNo{23}
%%%%%%%%%%%%%%%%%%%%%%%%%%%%%%%%%%%%%%%%%%%%%%%%%%%%%%

\begin{document}

\maketitle

\begin{abstract}
\end{abstract}

\section{Introduction}
\label{sec_intro}

\section{Background}
\label{sec_background}

\section{Preliminaries}
\label{sec_preliminaries}

%TODO: X introduce type var.
%TODO: X introduce formulas.
%TODO: X specify that formulas have type "var formula".
%TODO: X introduce sizef and maybe |.| notation.
%TODO: introduce the "set :: 'a list => 'a set" function.
%TODO: X introduce satisfiable
%TODO: X introduce the entailment relation
%TODO: X introduce formula equivalence
%TODO: X introduce disjunctive formulas
%TODO: X introduce tautological formulas
%TODO: introduce the "dualize :: 'a formula => 'a formula" function
A \emph{formula} is inductively defined as an atom, $\bot$, the negation of a formula, or the disjunction, conjunction or implication of two formulas. All atoms contained in such formulas are of the type \textsf{var}, representing a combination of a Boolean value and a natural number. A \emph{valuation} is defined as a function that maps an atom to a Boolean value. A valuation \emph{satisfies} a formula if the formula evaluates to True when all of its atoms are replaced by the specific truth values assigned to them by the valuation. A formula is \emph{satisfiable} if there exists a valuation that satisfies it. A formula is \emph{tautological} if it is satisfied by every valuation. Two formulas are \emph{equivalent} if they have the same truth value for every valuation. A \emph{literal} is a special case of a formula consisting of either a positive or negative atom, or a positive or negative $\bot$. A formula is \emph{conjunctive} if it is either a literal or its first subformula is a literal and its second a conjunctive formula. Analogously a formula is \emph{disjunctive} if it is either a literal or its first subformula is a literal and its second a disjunctive formula.

\begin{definition}The function set :: 'a list $\Rightarrow$ 'a set converts a list into a set.
\end{definition}

\begin{definition}The function sizef :: 'a formula $\Rightarrow$ nat calculates the size of a formula but without taking $\lnot$ into account.
\end{definition}

\begin{definition}The function dualize :: 'a formula $\Rightarrow$ 'a formula converts a formula in its dualized form by replacing each literal with its complementary literal and switching the symbols $\bot$ and $\top$, as well as $\land$ and $\lor$.
\end{definition}


\section{CNF to DNF}
\label{sec_CNF_to_DNF}

\begin{theorem}
  Let $n > 0$ be a natural number.
  Let $Ts$ be a list of formulas.
  Let $G = \bigvee_{T \in \mathsf{set} \, Ts} T$.
  If all formulas in $Ts$ are conjunctive and satisfiable (i.e., $\forall T \in \mathsf{set} \, Ts. \: \exists \alpha. \: \alpha \models T$), and $F_n$ is equivalent to $G$ (i.e., $\mathsf{equiv} \, F_n \, G$),
  then $\mathsf{sizef} \, G \ge n \cdot 2^n$.
\end{theorem}

\section{DNF to CNF}
\label{sec_DNF_to_CNF}

\begin{corollary}
  Let $n > 0$ be a natural number.
  Let $Cs$ be a list of formulas.
  Let $G = \bigwedge_{C \in \mathsf{set} \, Cs} C$.
  If all formulas in $Cs$ are disjunctive and not tautological (i.e., $\forall C \in \mathsf{set} \, Cs. \: \neg (\forall \alpha. \: \alpha \models C)$), and $\mathsf{dualize} \, F_n$ is equivalent to $G$ (i.e., $\mathsf{equiv} \, (\mathsf{dualize} \, F_n) \, G$),
  then $\mathsf{sizef} \, G \ge n \cdot 2^n$.
\end{corollary}

\section{Discussion}
\label{sec_discussion}

\section{Conclusion}
\label{sec_concl}

%%
%% Bibliography
%%

%% Please use bibtex,

\bibliography{bib.bib}

\end{document}
